% ═════════════════════════════════════════════════════════════════════════════
%  第六章  结论
% ═════════════════════════════════════════════════════════════════════════════
\section{结论}
\label{sec:conclusion}

\subsection{研究总结}

针对航空发动机装配领域知识密集、来源异构、查询精确性要求高的特点，
本文提出了面向 PT6A 涡桨发动机的多源知识图谱增强混合检索生成（MHRAG）框架。
全文工作可总结为四个层面。

\textbf{在知识构建层面}，本文设计了以 BOM 解析、CAD 结构化解析与
维修手册 LLM 抽取为三轨的多智能体协同流水线，
通过三级级联实体对齐机制将三类异构数据源汇入统一的 Neo4j 知识图谱，
构建出含 7\,333 节点 / 8\,752 边、覆盖 7 类实体与 9 类关系的航空发动机装配本体。
针对 CAD 几何邻接关系，本文进一步设计了基于轴对齐包围盒（AABB）间距计算的
物理路径与基于装配树兄弟节点推断的逻辑路径，使 \textit{adjacentTo} 关系
在不同部署环境下均可获得合理覆盖。

\textbf{在检索架构层面}，本文提出三路并行检索机制，
将密集向量（BGE-M3）、稀疏检索（BM25）与
基于 BGE-M3 语义近邻搜索的知识图谱子图召回路径
通过互惠排名融合（RRF）统一排序，
并经 BGE-Reranker 交叉编码器精排送入大语言模型完成生成。
其中 Dense 负责语义召回，BM25 负责精确字符串匹配，
KG 子图负责显式提供装配层级、工序依赖和技术规范关系。

\textbf{在生成与交互层面}，本文设计了多查询扩展、
"逐步确认"模式的结构化系统提示与 SSE 流式响应的端到端生成链路，
答案附带来源片段和相关图谱关系，便于前端展示证据来源。

\textbf{在系统验证层面}，本文构建了双层 QA 评测集
（Tier-1 精评 60 题 + Tier-2 跨章节粗评 20 题）。
MHRAG 在 Tier-1 生成阶段取得 Relevance 3.50、Completeness 3.02 与
Correctness 2.35 的表现，其中 Correctness 相对 Dense-only 提升 10.8\%；
在 Tier-2 跨章节场景下，相对 BM25-only 在 Completeness 与 Correctness 维度
分别提升 8.9\% 与 1.9\%。
工程性能维度，MHRAG 在 7.76\,s 端到端延迟和 1.60\,GB 本地显存占用下
取得 35.00\% 的答非所问率，体现出面向工业现场部署的可行性。

\subsection{研究创新点}

本文相对现有工作的主要贡献可归纳为三点：

\begin{enumerate}
  \item \textbf{三源知识图谱构建}。
        相比 Liu 等\cite{liu2023avkg}的 5 类实体 / 3 类关系航空装配 KG
        与 Shi 等\cite{shi2022arkg}仅含 CAD 属性数据的 ARKG，
        本文将 BOM、CAD 与维修手册三源数据通过统一本体（7 类实体 / 9 类核心关系）
        与三级级联对齐汇入同一图谱，构建覆盖更完整的装配知识底座。

  \item \textbf{语义近邻 KG 检索}。
        相比 GraphRAG\cite{edge2024graphrag}及现有工业 RAG 工作普遍采用的
        关键词字符串匹配 KG 检索方法，
        本文将 BGE-M3 编码同时作用于查询与图谱实体描述，
        使中英混合查询能直接命中语义近邻锚点，
        降低专业术语翻译和别名表达对图谱召回的影响。

  \item \textbf{RRF 异质融合}。
        相比 HybridRAG\cite{shen2025hybridrag}的线性加权融合方案，
        本文采用互惠排名融合（RRF）在不依赖人工权重调参的前提下
        将文本块与 KG 子图统一排序，
        对单路检索失败具有更好的鲁棒性。
\end{enumerate}

\subsection{工程价值与应用前景}

本文工作的工程价值体现在三个方面。
其一，本文实现了完整的端到端工程栈
（FastAPI + Neo4j + Qdrant + React），可为航空装配领域 RAG 系统部署提供参考。
其二，分阶段 KG 构建管道（Stage 1 BOM / Stage 2 Manual / Stage 3 CAD）
配合 HITL 审阅门控，使航空领域专家可在关键节点介入校验，
契合工业场景对知识可信度的严格要求。
其三，QA 评测集（80 题）与内部基线配置可作为后续研究的对比起点。

应用前景方面，本文方法可推广至以下场景：
（1）其他型号航空发动机的装配引导系统（如 CFM56、PW1000G）；
（2）相似精度要求的工业装配领域（如重型机床、轨道交通等）；
（3）作为 MES、PLM 等制造执行系统的智能化问答前端。

\subsection{未来工作展望}

未来研究可在以下方向深化：

\textbf{评测体系扩展。}
目前 Tier-1 黄金集来自单一压气机子系统，且存在模板化问题生成带来的评测偏置。
未来将通过领域专家盲测、多人独立标注以及真实生产场景采集，
构建覆盖整机各章节的领域级 RAG 评测基准，
并尝试与公开工业问答基准建立可比性。

\textbf{KG 增量更新与版本管理。}
当前 KG 构建为离线批处理流水线，未支持手册增补、零件号变更等增量场景。
未来将引入图谱差异化更新机制，
结合 \cite{amershi2014power} 中的人在回路审阅理念，
建立面向工程现场的 KG 版本管理与回退能力。

\textbf{推理深度增强。}
当前 KG 检索路径以一跳邻域扩展为主，对涉及多步推理的复杂查询
（如"安装 X 之前需要先完成的所有工序及其工具"）尚未充分发挥图结构优势。
未来将引入基于 GraphRAG 的多跳路径搜索与
Planner-Agent 迭代探索算法\cite{shen2025hybridrag}，
进一步提升 KG 路径在复杂结构化查询上的表现。
